\chapter{Nilpotent blocks and a common projector}\label{tx:sec:block-algebra}
\label{tx:ch:square-zero}
\section*{The property we need}
The lifting contract asks for one operation which works on an entire
proof diagram. It must transport the input-axiom map, the maps used for
multiplication, the inclusions of consequence spaces and the actual
target map. The first problem is to find an operation which has the
right effect on bases. The second, more difficult problem is to make
that operation exact on all the subspaces arising from a proof.

This chapter separates those problems. In characteristic two, a
square-zero operator gives a very concrete quotient: vectors killed
by the operator, modulo vectors in its image. On a permutation basis,
this quotient deletes the block labels. That calculation supplies
actual basis and tensor maps. It does \emph{not} prove exactness.
For exactness we will derive a projector from one group-algebra
witness, common to every tuple length up to the chosen support budget.
The next three chapters construct that witness rather than leaving
its existence as a hypothesis.

Keep the five roles separate:
\begin{center}
\begin{tabular}{@{}ll@{}}
$n$ & number of retained universe labels,\\
$q=2^a$ & number of block labels to be deleted,\\
$N=n+q$ & size of the larger universe,\\
$r$ & support budget for all spaces in the proof diagram,\\
$b\le r$ & current tuple length in the witness construction.
\end{tabular}
\end{center}
Neither $b$ nor $r$ is the maximum degree of a polynomial proof.
For an application to degree-$d$ consequences, that degree determines
$r$: $r=kd+s$ for invariant targets and $r=2kd+s+t$ for the general
indexed request. We first solve the finite algebra problem with $r$
as an independent integer.

The construction proves Exact Deletion
(Theorem~\ref{tx:thm:deletion}) independently of the equations and targets.
We begin with the square-zero operator in characteristic two,
then build a witness with a quadratic retained-label cost, replace that
cost by the linear threshold, and finally handle odd characteristic.

\section{A block of labels and a square-zero operator}\label{tx:sec:block2}
For this part $\operatorname{char}K=2$. Choose $q=2^a>r$ and split the labels into retained labels $R=[n]$ and a block $B$ with $q$ labels. Identify $B$ with the vector space $A=\Ftwo^a$. Translation by the $i$th unit vector gives a permutation $\sigma_i$ of $B$, extended by the identity on $R$. These permutations commute and square to the identity.

Extend the coefficient field if necessary so that it contains $a$
elements linearly independent over $\Ftwo$. In this finite-algebra
construction, $K$ denotes that working field; Lemma~\ref{tx:lem:field}
later returns degree-bounded derivability to the original field.
Choose $c_1,\ldots,c_a$ linearly independent over $\Ftwo$ in $K$, and put
\begin{equation}\label{tx:eq:z2}z=\sum_{i=1}^ac_i(\sigma_i-1).\end{equation}
This notation means an operator on every vector space on which label permutations act. Products of operators mean composition, so $z^2=0$ means that applying $z$ twice sends every vector to zero; it does not mean that $z$ itself is zero. Since the summands commute, cross terms cancel in characteristic two, and
\[z^2=\sum_i c_i^2(\sigma_i-1)^2=0.\]
A linear operator is called \emph{nilpotent} when some positive
power is zero. Here that power is already two, so every vector
in $\im z$ lies in $\ker z$. The nonzero two-point example below
shows that nilpotence need not make the operator itself zero.
For $q=2$, $z$ can be the swap of the two block labels plus the identity. For $q=4$, work over $\mathbb F_4=\Ftwo[\omega]/(\omega^2+\omega+1)$ and use $(c_1,c_2)=(1,\omega)$.

The independence is essential. On a two-point orbit where $\sigma_1$ and $\sigma_2$ induce the same swap, choosing $c_1=c_2=1$ makes $z=0$. The orbit is moved, yet would survive deletion. Independent coefficients ensure that every nonzero binary combination of the $c_i$ remains nonzero.

For a space $M$ with label action define
\begin{equation}\label{tx:eq:D2}\D(M)=\ker z/\im z.\end{equation}
This quotient is legal because $z^2=0$ implies $\im z\subseteq\ker z$. Two vectors killed by $z$ represent the same class if their difference is $zw$ for some $w$. A map commuting with label permutations commutes with $z$ and induces a map between these quotients. Permutations of the retained labels commute with $z$, so they still act on the quotient.

\subsection{Two actions which must not be confused}
The coordinates of a tuple have \emph{positions}, whereas its entries
are \emph{universe labels}. For example, the tuple $(a,b)$ has two
positions and uses two labels. A label permutation $g$ sends it to
$(g(a),g(b))$. A position swap sends it to $(b,a)$. These actions
commute, but they play different roles. Group-algebra elements in
$K[S_N]$ combine label permutations. The auxiliary inverses used later
must commute with position permutations so that their matrix entries
can be expressed as functions of an unordered matching between input
and output labels.

For a concrete check, take labels $1,2,3$ and the tuple $(1,3)$.
The label swap $(12)$ sends it to $(2,3)$; the position swap sends
it to $(3,1)$. Doing both, in either order, gives $(3,2)$. It is
this commuting pair of actions which allows the same correction to
be discussed both as an operator on ordered tuples and as a function
on rook placements.

The block subgroup $A$ acts on labels, never on tuple positions. The
operator $z$ therefore commutes with every position permutation.
An auxiliary operator chosen by the underlying \emph{set} of labels,
instead of by a preferred ordering of a tuple, has the same property.
That observation will matter when constructing an inverse without
averaging over $S_b$.

The usual averaging recipe would replace an operator $h$ by
$(b!)^{-1}\sum_{\tau\in S_b}\tau h\tau^{-1}$, thereby forcing
it to commute with positions. In characteristic two, $b!$ is zero
in the field when $b\ge2$, so this recipe is unavailable. Instead
we construct $h$ from the unordered label set of a tuple. That
choice is already unchanged by every reordering of its positions.

\subsection{The two-point orbit calculation}
For the illustrative two-label choice $z=\sigma-I$, on basis vectors $e_\beta,e_\gamma$ swapped by $\sigma$,
\[ze_\beta=ze_\gamma=e_\beta+e_\gamma.\]
The kernel and image are both the one-dimensional span of $e_\beta+e_\gamma$, so the quotient is zero. On a fixed basis vector $e_i$, $z=0$, and its line survives. This is the prototype of deletion: moved directions disappear, fixed directions remain.

\paragraph*{A useful failure to commute.}
On this two-point orbit let $S$ be the swap and let $D$ multiply
the first basis vector by $0$ and the second by $1$. Their matrices
are
\[
 S=\begin{pmatrix}0&1\\1&0\end{pmatrix},\qquad
 D=\begin{pmatrix}0&0\\0&1\end{pmatrix},\qquad DS-SD=S
 \quad(\operatorname{char}K=2).
\]
The difference $DS-SD$ is called a \emph{commutator}; it measures
the failure of the two operators to commute. Here that difference
is invertible. The next calculation constructs the same useful
failure to commute on every moved orbit, and uses it to write any
kernel vector as an image vector.

\subsection{The calculation for every moved orbit}
An $A$-set is a set of basis objects on which the translations in
$A=\mathbb F_2^a$ act. The orbit of $x_0$ consists of all $v x_0$,
$v\in A$; its stabiliser $H$ consists of translations fixing $x_0$.
Since the orbit is moved, $H$ is a proper subspace. A basis of $H$
can be extended to a basis of $A$, so a nonzero linear functional
$\ell$ can be chosen to vanish on $H$. It assigns a binary coordinate
to the orbit, independent of the translation used to reach an object.
The coordinate-reading operator $D$ below is diagonal, just as in
the two-point illustration.

Let $X$ be a finite $A$-set with no point fixed by all of $A$. On each orbit choose $x_0$. Its stabiliser $H$ is a proper binary subspace. Choose a nonzero linear functional $\ell:A\to\Ftwo$ vanishing on $H$. The function
\[\pi(vx_0)=\ell(v)\]
is well defined: two choices of $v$ differ by an element of $H$. Let $D$ multiply $e_x$ by $\pi(x)$, and put $e_i=\ell(\varepsilon_i)$. These are constants on that orbit. If $S_i$ denotes the action of $\sigma_i$, direct evaluation at a basis vector gives
\[DS_i-S_iD=e_iS_i,\qquad Dz-zD=W:=\sum_i c_ie_iS_i.\]
Moreover $Wz=zW$ and
\[W^2=\Bigl(\sum_i c_ie_i\Bigr)^2I=\alpha^2I,\qquad\alpha\ne0.\]
In particular $W^{-1}=\alpha^{-2}W$. Thus the inverse used here
is explicit; no choice of a complement or inverse matrix is hidden
in the notation. The coordinate vector $\varepsilon_i$ is the
$i$th unit vector in $A$, and $e_i=\ell(\varepsilon_i)$ is a
scalar, distinct from a basis vector $e_x$.
Thus $\Theta=W^{-1}D$ satisfies $\Theta z-z\Theta=I$. For $zm=0$ this gives $m=-z\Theta m$, hence $\ker z\subseteq\im z$. The reverse inclusion was automatic, so the quotient vanishes.

The calculation has three ingredients, each serving a different
purpose. The orbit coordinate $\pi$ produces a diagonal operator whose
commutator with translation is computable. Independence of the $c_i$
makes the commutator invertible. The commutator identity then turns a
kernel vector into an image vector. We do not choose a complement of
$\ker z$ by hand and hope it will respect later maps.

For a two-point orbit of $A=\mathbb F_2^2$ with stabiliser
$H=\{(0,0),(1,1)\}$, both coordinate translations induce the same
swap $S$. Choose $\ell(v_1,v_2)=v_1+v_2$. With
$(c_1,c_2)=(1,\omega)$, the operator is
$z=(1+\omega)(S-I)$, and $\alpha=1+\omega\ne0$.
Thus the one-dimensional kernel is exactly the image. If both
coefficients had instead been $1$, $z$ would be zero on this moved
orbit. This explicit orbit is why using one coefficient for every
translation is not a harmless simplification.

This calculation also works on $K[X]\otimes Y$ for \emph{any} space $Y$ carrying the label action. Let $D$ inspect only the $X$ coordinate and let $S_i$ act on both factors. We still have $S_i^2=I$ and the same commutator calculation. This stronger statement will justify multiplication, whose domain contains a variable factor tensored with a consequence subspace.

\section{Natural deletion on permutation bases}
Let $\mathcal U$ be a basis system whose objects have support less than $q$. An object fixed by all of $A$ cannot have support meeting $B$. Otherwise invariance of its support under the regular action would force its support to contain all $q$ block labels. Thus its fixed objects are precisely $\mathcal U_n$.

Split the basis into fixed and moved objects. The preceding calculation kills the moved summand and fixes the retained one. Hence
\begin{align}
\D(K[\mathcal U_{n+q}])&\cong K[\mathcal U_n],\label{tx:eq:basis}\\
\D(K[\mathcal U_{n+q}]\otimes Y)&\cong K[\mathcal U_n]\otimes\D(Y).\label{tx:eq:tensor}
\end{align}
The second identification sends $e_u\otimes[y]$ to $[e_u\otimes y]$ for retained $u$. It is natural in $Y$: applying a map to $y$ before or after taking its class gives the same result. For a product of two permutation bases, it simply keeps pairs of retained objects. There is no choice of unrelated isomorphisms here.

Suppose a matrix from one permutation basis to another has stable relabelling-invariant coefficients. On a retained input vector, its image is killed by $z$. Its retained component has exactly the smaller-rank coefficients; its moved component is in $\ker z=\im z$ on the moved summand, so vanishes in the quotient. Thus the \emph{actual matrix}, not merely its rank, is transported correctly.

\subsection{What naturality says about actual matrix entries}
Let $F_N:K[\mathcal U_N]\to K[\mathcal V_N]$ be an equivariant
coefficient map. On a retained basis vector $e_u$, its image is
\[
                       F_N(e_u)=\sum_v a_N(v,u)e_v.
\]
It is killed by $z$ because $e_u$ is killed by $z$ and $F_Nz=zF_N$.
Split this image into its retained and moved parts. Both parts are
killed separately because $z$ preserves those two summands. The moved
part is in $\im z$ by the orbit calculation, and so contributes zero
to the quotient. The remaining coefficients are $a_N(v,u)$ for
retained $u,v$. A stable coefficient rule identifies each of these
with $a_n(v,u)$. The smaller matrix is therefore obtained entry by
entry; we have not merely proved equality of two ranks or dimensions.

For a tensor map, apply the same argument to the permutation first
factor. A retained basis vector in that factor is fixed by the block,
so the quotient sees precisely the class of the second factor. A moved
first-factor orbit has zero quotient even when the block acts
nontrivially on the second factor. This is why the strengthened orbit
calculation, with a diagonal action on $K[X]\otimes Y$, was needed.
The tensor rule subsequently applies to a variable factor tensored
with a consequence subspace.

\subsection{A small failure of exactness}
Take the two-point permutation space $E=Ke_\beta\oplus Ke_\gamma$ and its invariant line $T=K(e_\beta+e_\gamma)$. Inclusion $T\hookrightarrow E$ is injective. But $\D(T)=T$ whereas $\D(E)=0$. Deletion takes an injection to the zero map on a nonzero space. Therefore the basis calculation alone does not prove that deletion preserves kernels, images or intersections.

This failure is precisely why a proof cannot stop at ``the monomials delete correctly.'' We must show that, with enough retained labels, every subspace produced by the PC calculation lies in a class where the quotient operation is exact.

\section{A common group-algebra witness}\label{tx:sec:common}
The group algebra $K[S_N]$ consists of finite formal sums $\sum_ga_gg$ of label permutations. It acts on every space with an $S_N$ action by the corresponding linear combination of operators. Its multiplication is composition. This provides a common language for certificates valid on many spaces at once.

For $b\ge0$ let
\[P_b(N)=K[\Inj([b],[N])],\]
the space on ordered injective $b$-tuples. In these construction chapters $P_b(N)$ denotes this tuple space, not the polynomial ambient space $P_{N,\le d}$ from the consequence calculation. The integer $b$ counts tuple positions, whereas $d$ bounds polynomial degree. Label permutations act on entries; a second group $S_b$ permutes tuple positions. These actions commute. Define
\[J_r(N)=\{a\in K[S_N]:aP_b(N)=0\text{ for every }b\le r\}.\]
It is a two-sided ideal. Write $\mathcal C_r(N)$ for the finite-dimensional spaces with label action annihilated by this ideal: every element of $J_r(N)$ acts as zero on each such space. This is the same annihilator class called $\mathcal M_r$ in Chapter~\ref{tx:chap:lift}. Subspaces invariant under labels, quotients and finite direct sums remain in this class.

Every bounded-support permutation basis of support at most $r$ belongs to the class. For an orbit of an object, order its support of size $b\le r$. Sending an injective tuple to the corresponding relabelling of that object gives a surjective equivariant map from $P_b(N)$ onto the orbit space. The support condition makes this well defined. Elements of $J_r$ kill that quotient as well.

The required witness is a \emph{partial inverse}, not an inverse of
the nilpotent operator $z$. If $zwz=z$ and $v=zt$ is an image
vector, then $z(wv)=zwzt=zt=v$: applying $w$ supplies a preimage
under $z$. The witness condition below requires this one expression
to work on all tuple lengths simultaneously. Its error need only
lie in $J_r$, since such an error acts as zero on the entire class.

\begin{definition}
A support-$r$ witness is one element $w\in K[S_N]$ satisfying
\begin{equation}\label{tx:eq:witness2}zwz-z\in J_r(N).
\end{equation}
\end{definition}
Such a $w$ gives $zwz=z$ on \emph{every} object of $\mathcal C_r(N)$. Choosing a separate matrix inverse on each space would not give this compatibility. A module map here is a linear map commuting with all label permutations. It therefore satisfies $f(wv)=w f(v)$ for a common group-algebra element, which is exactly what is needed.

\subsection{The projector calculation, without an exactness assumption}
A projector is an operator $e$ with $e^2=e$. It fixes every vector
in its image $eM$ and sends every vector in $\ker e$ to zero;
it need not be a coordinate projection.
We seek one which selects representatives
of the quotient $\ker z/\im z$ in a way compatible with every
equivariant map. The factors in the formula below act from right
to left. They are not assumed to commute.

Suppose $z^2=0$ and $zwz=z$. Set
\[e=(I-zw)(I-wz).\]
Then
\[ze=z(I-wz)=z-zwz=0,\qquad ez=(I-zw)z=0.\]
If $zv=0$, then $ev=v-zwv$. Also
\[(zw)e=(zw-zwzw)(I-wz)=0,\]
since $zwzw=zw$. Applying the formula on $\ker z$ to $ev$ gives $e^2v=ev$. Thus $e$ is a projector.

The map $[v]\mapsto ev$ is an isomorphism $\ker z/\im z\to eM$: it is well defined because $ez=0$, its image lies in $\ker z$, and $ev-v=-zwv$ is in $\im z$ when $zv=0$. Its inverse sends a vector of $eM$ to its quotient class.

Now let $0\to U\to V\to W\to0$ be an exact sequence in $\mathcal C_r(N)$. All maps commute with the same $e$. Restricting to $eU,eV,eW$ preserves injectivity. A vector of $eW$ has a lift $v\in V$; $ev$ is a lift in $eV$, proving surjectivity. A vector $ev$ mapping to zero lies in $U$ and equals its own projection, so lies in $eU$. This proves exactness explicitly. By the natural isomorphism, $\D$ is exact.

The projector need not commute with permutations of the retained labels. We use its image as a vector-space model to prove exactness, and keep the original quotient for the residual label action. Nothing here asserts an unwarranted equivariant direct-sum decomposition for that residual group.

\begin{example}[One retained label]
Let the retained label be $a$ and the two block labels be $b,c$. Put $z=I+(bc)$ and $w=(ab)$. On $P_1(3)$,
\[zwz(e_b)=zw(e_b+e_c)=z(e_a+e_c)=e_b+e_c=ze_b,\]
and similarly for $e_c$; both sides vanish on $e_a$. The projector is the all-ones $3\times3$ matrix. Its image is the constant line, naturally representing the retained one-dimensional quotient. With no retained label, every element of $K[S_2]$ satisfies $zwz=0$, so there is no such witness. The spare label makes a real difference.
\end{example}


\section{Reading the projector as an exactness certificate}
In a short exact sequence the first map $i$ is injective, the last
map $f$ is onto, and $\ker f=\im i$. Thus $U$ is a specified
subspace of $V$, and $W$ records the quotient by it. A
\emph{diagram chase} checks these assertions after transport by
following an individual vector and its preimages through the maps.
Here the common projector makes those vector checks possible.

The exactness proof is worth restating as a diagram chase, because
it is the point at which a basis calculation becomes a tool for
ordinary PC. Suppose
\[
 0\longrightarrow U\xrightarrow{i}V\xrightarrow{f}W\longrightarrow0
\]
is exact and all three spaces are in the common annihilator class.
The equations $ei=ie$ and $ef=fe$ hold because $i,f$ commute with
every label permutation and hence with their common linear combination
$e$. A separate projector selected independently on each space would
not justify these equations.

If $u\in eU$ maps to zero, its original image is zero and injectivity
of $i$ gives $u=0$. If $w\in eW$, take any lift $v\in V$ with
$f(v)=w$. Then $f(ev)=e f(v)=ew=w$. Finally, if $v\in eV$ maps
to zero, exactness of the original sequence gives $v=i(u)$ for some
$u\in U$. Applying $e$ gives $v=ev=i(eu)$. These are respectively
injectivity, surjectivity and equality of image with kernel for the
projected sequence. The natural quotient-to-projector isomorphism
transfers all three conclusions to deletion.

This argument proves more than the existence of some complement to
$\im z$. It proves that every actual map in the consequence diagram
uses compatible complements. Intersections, sums and restricted
multiplication images can therefore be transported. The joint tensor
operation in \eqref{tx:eq:joint-lift} uses the same group-algebra
expression on both factors together; it is not the tensor product
of two independently selected projectors.

\subsection{The three-label matrix in full}
With retained label $a$ and block labels $b,c$, use the ordered basis
$(e_a,e_b,e_c)$. Then
\[
 z=\begin{pmatrix}0&0&0\\0&1&1\\0&1&1\end{pmatrix},\qquad
 w=\begin{pmatrix}0&1&0\\1&0&0\\0&0&1\end{pmatrix}.
\]
Matrix multiplication gives $zwz=z$ and
\[
 e=(I-zw)(I-wz)=
       \begin{pmatrix}1&1&1\\1&1&1\\1&1&1\end{pmatrix}.
\]
Its square is itself: every entry of the square is $1+1+1=1$ in
characteristic two. The quotient can be represented by the retained
coordinate, while the projector realises that coordinate as
$e_a+e_b+e_c$. Thus the lift of a retained vector can involve all
labels. This is already the mechanism that turns a smaller totality
sum into its larger counterpart.

The example also shows why deletion is not substitution of zero for
block variables. The projector lift of $e_a$ is a sum of three basis
vectors. Its block terms are required to preserve the relations of
the larger system. A support budget does not assert that lifted proof
lines mention only retained labels; it limits the labels in each
\emph{basis object} used by the transport construction.

\section*{What has been proved, and what remains}
We have proved natural deletion of permutation bases and its
permutation-first-factor tensor rule. We have also proved that the
single identity $zwz=z$ on all tuple lengths up to $r$ gives a common
projector and therefore exact deletion on the full annihilator class.
The identity has not been assumed as an external theorem: its
construction is the subject of the next chapters. The preliminary
construction spends fresh labels at each length; the improved
construction realises a smaller, specially shaped error on one reused
pool.
